\section{СРАВНИТЕЛЬНЫЙ МЕТОД В ПОЛИТИЧЕСКОЙ НАУКЕ: СУЩНОСТЬ И ПРОБЛЕМЫ}

\subsection{Сравнение как метод и субститут эксперимента}
Сравнение выступает фундаментальной установкой научного познания. В политической науке сравнительный метод занимает особое место, поскольку проведение контролируемого научного эксперимента в макрополитической сфере практически невозможно. В силу этого, как отмечал А. Лейпхарт, сравнение выступает наиболее адекватным заменителем (субститутом) экспериментального метода~\cite{smorgunov2012}.

Как инструмент анализа сравнение позволяет обнаруживать корреляционные и каузальные зависимости между политическими феноменами. Научное признание метода восходит к концепции Дж. С. Милля, который указывал на необходимость верификации эмпирических законов общества через серийные исторические наблюдения и их соотнесение с общими законами человеческой природы~\cite{smorgunov2012}. Таким образом, компаративный метод выступает важнейшим связующим звеном между абстрактным политико-философским знанием и конкретными эмпирическими событиями~\cite{smorgunov2012}.

\subsection{Соотношение теории и метода в компаративистике}
Взаимосвязь теории и метода в сравнительной политологии имеет диалектический характер. С одной стороны, сама дисциплина получила свое название по методу, а не по предмету, что нередко приводило к дискуссиям о растворении сравнительной политологии в общем поле политической науки~\cite{smorgunov2012}. С другой стороны, сравнение коренным образом трансформирует сам теоретический результат исследования, придавая получаемым обобщениям синтезированный характер~\cite{smorgunov2012}.

В современной политологии теория перестает играть лишь чисто инструментальную роль классификатора эмпирических данных. Она становится конечной целью сравнительного анализа~\cite{smorgunov2012}. Теоретические модели (например, концепции рационального выбора или системные абстракции) задают первоначальные критерии релевантности переменных, определяют рамки конструирования гипотез и защищают исследователя от ложных корреляций. В то же время именно сравнительный метод обеспечивает верификацию и фальсификацию выдвинутых макротеоретических положений на основе кросс-национальных или кросс-временных данных.